%!TEX root = ../main/main.tex

\section*{Instrucciones}

Cualquier duda sobre la tarea se deberá hacer en los \emph{issues} del \href{https://github.com/IIC3253/2026}{repositorio del curso}. Los issues son el canal de comunicación oficial para todas las tareas.

\paragraph{Configuración inicial.} 
Para esta tarea utilizaremos \textit{github classroom}. 
Para acceder a su repositorio privado debe ingresar a un link que le enviaremos antes de la entrega de la tarea, seleccionar su nombre y aceptar la invitación.
El repositorio se creará automáticamente una vez que haga esto y lo podrá encontrar junto a los repositorios del curso.
Para la corrección se utilizará Python 3.13.7.

\paragraph{Entrega.} Al entregar esta tarea, su repositorio se deberá ver exactamente de la siguiente forma:

\bigskip

\dirtree{%
.1 \faGithub \ Repositorio.
.2 \faFolderOpenO \ \texttt{Tarea2}.
.3 \faFilePdfO \ \texttt{feistel.py}.
.2 \faFileTextO \ \texttt{.gitignore}.
.2 \faFileTextO \ \texttt{README.md}.
.2 \faFolderO \ .git.
}

